\documentclass[11pt]{article}
\usepackage[margin=1in]{geometry}
\usepackage{amsmath,amssymb,amsthm}
\usepackage[hidelinks]{hyperref}
\newtheorem{theorem}{Theorem}
\newtheorem{proposition}[theorem]{Proposition}
\newtheorem{corollary}[theorem]{Corollary}
\newtheorem{remark}[theorem]{Remark}
\newcommand{\Prb}{\mathbb P}
\newcommand{\Bzero}{\mathcal B_0}
\newcommand{\Bone}{\mathcal B_1}
\newcommand{\errzero}{\mathcal E_0}
\newcommand{\errone}{\mathcal E_1}
\title{Truth-uniform error contracts for monotone queries with shared sources}
\author{Alexander Richard Towell\\Southern Illinois University Edwardsville\\
\texttt{lex@metafunctor.com}}
\date{October 13, 2026\\\small Internal research note; originality assessment incomplete}

\begin{document}
\maketitle
\begin{abstract}
Approximate-query error propagation usually conditions on a fixed exact
database or assumes a distribution over exact databases. We consider instead
the worst error over all exact truth assignments giving the specified answer.
For a monotone Boolean lineage and independent one-sided source noise of rate
$\alpha$, the exact worst false-positive probability is
$1-(1-\alpha)^{s_0}$, where $s_0$ is the function's 0-sensitivity. The dual
false-negative formula uses 1-sensitivity. Heterogeneous rates give weighted
optimization over boundary certificates; arbitrary dependence gives the sharp
homogeneous envelope $\min\{1,\alpha s_0\}$. Shared tuple occurrences must first
be identified as single sources. On binary-witness lineage of source occurrence
at most two, the sensitivity has a linear-time graph formula. These elementary
contract consequences are proved and exhaustively checked on small functions.
They are not asserted to be a new general probability calculus or a sufficient
standalone publication contribution.
\end{abstract}

\section{Question, semantics, and scope}

Fix a monotone function $F:\{0,1\}^n\to\{0,1\}$ describing the Boolean
lineage of one query answer. The variables index distinct frozen random
sources, not syntactic tuple occurrences. A canonical constraint such as
$R(u,v)=R(v,u)$ identifies those two tuple occurrences with one variable.
Repeated evaluation of that source reads the same bit. Distinct sources are
independent only where explicitly assumed.

Let $x$ be the exact source-membership vector. We permit every $x\in\{0,1\}^n$.
This is appropriate for class-induced constraints after quotienting because
every assignment to the classes defines an admissible exact relation. Additional
constraints, such as functionality or transitivity, can restrict admissible
truth vectors. The formulas below then remain upper bounds but need not be sharp.

Under independent one-sided false-positive noise with rates
$\boldsymbol\alpha=(\alpha_i)_{i=1}^n$, use independent uniforms $U_i$ and set
\[
 Y_i^+(x)=x_i\lor\mathbf 1\{U_i\leq\alpha_i\}.
\]
Under independent one-sided false-negative noise with rates
$\boldsymbol\beta$, set
\[
 Y_i^-(x)=x_i\land\mathbf 1\{U_i>\beta_i\}.
\]
Define the truth-uniform errors
\begin{align*}
 \errzero(F;\boldsymbol\alpha)
 &=\max_{x:F(x)=0}\Prb\{F(Y^+(x))=1\},\\
 \errone(F;\boldsymbol\beta)
 &=\max_{x:F(x)=1}\Prb\{F(Y^-(x))=0\}.
\end{align*}
If the relevant class of exact inputs is empty, define the corresponding
error as zero. Constant functions therefore have both errors zero.
This is a mathematical convention for the extremal formulas; the companion
prototype returns \texttt{None} when the conditioning class is empty, because
the corresponding conditional error rate is undefined.
No prior on $x$ is assumed. This differs from computing the probability of a
fixed lineage under one known source-probability vector.

Let $\Bzero$ be the inclusion-maximal false vectors and $\Bone$ the
inclusion-minimal true vectors. The one-sided sensitivities are
\[
 s_b(F)=\max_{x:F(x)=b}
 \bigl|\{i:F(x\mathbin\oplus e_i)\ne F(x)\}\bigr|,
 \qquad b\in\{0,1\},
\]
with value zero when the corresponding input class is empty.

\section{Exact independent-source contracts}

\begin{theorem}[Boundary formula]
For any monotone nonconstant $F$,
\begin{align}
 \errzero(F;\boldsymbol\alpha)
 &=\max_{b\in\Bzero}
 \left[1-\prod_{i:b_i=0}(1-\alpha_i)\right],\label{eq:hetero-fp}\\
 \errone(F;\boldsymbol\beta)
 &=\max_{b\in\Bone}
 \left[1-\prod_{i:b_i=1}(1-\beta_i)\right].\label{eq:hetero-fn}
\end{align}
\end{theorem}

\begin{proof}
If $x\leq x'$, the uniform construction gives
$Y^+(x)\leq Y^+(x')$ pointwise. Monotonicity of $F$ implies that its acceptance
probability is nondecreasing in the exact profile. Every false $x$ has an
extension $b\in\Bzero$, so it suffices to maximize over such $b$.
At a maximal false $b$, turning any remaining zero into a one makes $F$ true.
Existing ones never disappear. Hence $F(Y^+(b))=0$ if and only if all zeros
of $b$ stay zero. Independence proves~\eqref{eq:hetero-fp}.

For the false-negative case, shrink every true input to a minimal true input
and use the same monotone coupling. At a minimal true input, losing any
remaining one makes $F$ false, and no zero can become one. Independence proves
\eqref{eq:hetero-fn}. Both maxima are attained on actual exact profiles.
\end{proof}

\begin{proposition}[Boundary sizes]
For monotone nonconstant $F$,
\[
 s_0(F)=\max_{b\in\Bzero}|\{i:b_i=0\}|,
 \qquad
 s_1(F)=\max_{b\in\Bone}|\{i:b_i=1\}|.
\]
These are also the standard one-sided certificate complexities.
\end{proposition}

\begin{proof}
Every zero at a maximal false vector is pivotal, yielding one inequality.
Conversely, extend any false vector $x$ to a maximal false vector $b$.
A pivotal zero of $x$ cannot have been changed to one during that extension:
otherwise monotonicity would force $F(b)=1$. It remains pivotal at $b$,
again by monotonicity. Thus no pivotal zero is lost. The true-input case is
dual. For the certificate statement, the zeros of a maximal false vector
form a minimal zero certificate, and every false vector has a zero certificate
arising from some maximal false extension. The true case is analogous.
\end{proof}

\begin{corollary}[Homogeneous contract]
When all false-positive rates equal $\alpha$ and all false-negative rates
equal $\beta$ in their respective one-sided models,
\[
 \boxed{\errzero(F;\alpha)=1-(1-\alpha)^{s_0(F)}},\qquad
 \boxed{\errone(F;\beta)=1-(1-\beta)^{s_1(F)}}.
\]
\end{corollary}

\begin{remark}[Two-sided channels]
The two expressions remain upper bounds when both error types are possible:
removing false negatives can only increase false positives, and removing
false positives can only increase false negatives. They need not be sharp.
For $F=X_1\land X_2$, homogeneous rates with $\alpha\leq1-\beta$ give exact
worst false-positive probability $\alpha(1-\beta)$, whereas the one-sided
formula is $\alpha$. False positives can rescue false negatives, so these
effects cannot be combined by simply declaring both one-sided identities exact.
\end{remark}

\section{Sharp contracts without independence}

Now allow arbitrary dependence among the one-sided error events, subject
only to the per-source marginal upper bounds. The supremum is over admissible
couplings as well as exact profiles.

\begin{theorem}[Dependence-robust contract]
With homogeneous marginal false-positive bound $\alpha$,
\[
 \sup_{x:F(x)=0,\ \text{couplings}}
 \Prb\{F(Y^+(x))=1\}
 =\min\{1,\alpha s_0(F)\}.
\]
For heterogeneous marginal bounds the sharp expression is
\[
 \max_{b\in\Bzero}
 \min\left\{1,\sum_{i:b_i=0}\alpha_i\right\}.
\]
The false-negative statements are dual, using $\Bone$ and $s_1$.
\end{theorem}

\begin{proof}
For any false $x$, extend it to a maximal false $b$. Its remaining zero
coordinates form a zero certificate already present in $x$. If none of
these coordinates flips, the output stays false, regardless of all other
coordinates. The union bound proves the proposed upper bound.

At a maximizing maximal false profile, every zero flip is sufficient to
change the output. Events with prescribed probabilities $a_1,\ldots,a_k$
can have union probability $\min\{1,\sum_i a_i\}$: put disjoint intervals
of those lengths in the unit interval if their sum is at most one; otherwise
lay intervals consecutively on the unit circle, whose union covers the
circle. This supplies an attaining coupling for the source errors. The
dual argument handles false negatives.
\end{proof}

Source identity therefore changes the Boolean function and its certificate
invariant. Independence is separate information that replaces a union bound
by an exact product at the extremal certificate. For example, repeating the
same source $k$ times in an OR leaves $s_0=1$ and error $\alpha$;
$k$ genuinely distinct sources in an OR have $s_0=k$.

\section{Binary witnesses and practical closed forms}

Write positive query lineage as
\[
 F=\bigvee_{j=1}^m\bigwedge_{i\in S_j}X_i.
\]
Normalize the supports by identifying source aliases, removing repeated
sources inside a witness, deleting duplicate witnesses, and absorbing every
witness containing another. Then $s_1$ is the maximum witness size, while
$s_0$ is the largest size of an inclusion-minimal hitting set of the
witness family. The word ``minimal'' refers to inclusion, not cardinality.

For binary witnesses, represent every support $\{u,v\}$ by an edge of a
simple graph $G$. Maximal false profiles are maximal independent sets, hence
\begin{equation}
 s_0(F_G)=|V(G)|-i(G),\label{eq:independent-domination}
\end{equation}
where $i(G)$ is the minimum size of a maximal independent set, equivalently
the independent domination number. If an edge exists, $s_1(F_G)=2$.

\begin{proposition}[Occurrence at most two]
Suppose normalized witness width and source occurrence are at most two.
Let $U$ be the singleton witnesses, and let the remaining binary graph
have path or cycle components $G_c$ of orders $n_c$. Then
\[
 s_0=|U|+\sum_c\left(n_c-\left\lceil n_c/3\right\rceil\right),
 \qquad s_1=\max_j|S_j|.
\]
These quantities and the homogeneous one-sided contracts are computable
with a linear number of graph operations.
\end{proposition}

\begin{proof}
The occurrence bound makes graph degree at most two. A selected vertex in
an independent dominating set of a path or cycle covers at most three
vertices, so $i(G_c)\geq\lceil n_c/3\rceil$.
For a path, select every third vertex starting at the second when
$n_c\equiv0,2\pmod3$, and starting at the first when $n_c\equiv1\pmod3$.
For a cycle, let $r=\lceil n_c/3\rceil$ and put selected vertices at cyclic
gaps two or three summing to $n_c$; this is possible because
$2r\leq n_c\leq3r$. These sets are independent and dominating.
The invariant adds over components, and every singleton must belong to
every hitting set. Equation~\eqref{eq:independent-domination} completes the proof.
\end{proof}

Heterogeneous false-positive rates require weighted optimization over
minimal vertex covers. On a path, minimize
$\prod_{i\notin S}(1-\alpha_i)$ over independent dominating sets $S$.
A three-state dynamic program tracks whether the current vertex is selected,
unselected and dominated from the left, or unselected and awaiting a selected
right neighbor. With $q_i=1-\alpha_i$, its state update is
\[
 A'=\min(B,C),\qquad B'=q_iA,\qquad C'=q_iB.
\]
Initialize $(A,B,C)=(1,\text{unreachable},q_1)$ and accept states $A,B$
at the end. Unreachable states are not ordinary numerical infinities
multiplied by zero. Cycles require a constant number of boundary cases.
This is a standard independent-domination dynamic program with multiplicative
costs. The homogeneous formula avoids even this optimization.

\section{An adjacent limitation: symmetry and exact inference}

Identifying symmetric tuples preserves marginal rates, but can change the
complexity of exact inference. Let $G$ be a simple cubic graph, and let an
exact relation $D$ contain one tuple $(v,a,b,c)$ listing the three neighbors
of each vertex $v$. Consider the fixed query
\[
 \exists v,a,b,c:\quad D(v,a,b,c)\land R(v,a)\land R(v,b)\land R(v,c).
\]
Let all queried tuples be exact nonmembers, and let their FPR be $p$.
Independent directed tuple sources give read-once lineage and false-positive
probability $1-(1-p^3)^{|V|}$. Canonical symmetry identifies $R(u,v)$ with
$R(v,u)$; each undirected edge is now a single independent bit used twice.
Every individual tuple rate $p$ and witness rate $p^3$ is preserved.

The query is false exactly when the zero edges form an edge cover. At
$p=1/2$ its FPR is therefore
\[
 1-\frac{\#\operatorname{EdgeCovers}(G)}{2^{|E|}}.
\]
Counting edge covers on simple planar cubic graphs is already
\#P-hard~\cite{amarilliMonet,caiLuXia}. This is a consequence of established
counting results, not a new hardness theorem. It explains why a compact
worst-case contract and an exact probability for a particular truth profile
can be very different computational goals. A nonempty exact relation can
be supplied by adding an unqueried true symmetric sentinel pair, if the
underlying formalism treats the empty relation as a special exact case.

\section{Verification and novelty assessment}

The companion program \texttt{check\_contracts.py} enumerates every monotone
function on one through four variables: respectively 3, 6, 20, and 168
functions. It checks the homogeneous identities at rates
$0,1/5,1/2,1$, plus heterogeneous identities, using exact rational arithmetic.
All 1,970 identity checks passed. It also verified the graph invariant on
305 labeled maximum-degree-two graphs with at most five vertices.
The separate program \texttt{check\_theory.py} checked 256 projection-covariance
cases, eight weighted paths, 87 weighted cycles including degenerate rates,
and the cubic $K_4$ edge-cover construction. The latter has 41 edge covers,
canonical FPR $23/64$, and prequotient FPR $1695/4096$.

These checks validate implementations on a finite domain; the proofs, not
the tests, justify the general identities. No implemented approximate filter
has been shown to satisfy the independent-source model in this note.

Boolean provenance, weighted model counting, and dissociation already provide
sound ways to evaluate or bound a fixed truth profile~\cite{gatterbauerSuciu}.
Correlated databases can make even read-once lineage difficult~\cite{kanagalDeshpande}.
For monotone functions, sensitivity and certificate complexity coincide;
this is standard and explicitly recalled in recent work~\cite{blanc}.
Gopalan et al. give pointwise two-sided noise bounds in terms of sensitivity
\cite[Lemma 15]{gopalan}. The present sharp one-sided extremal formulas are
elementary certificate consequences. The targeted search did not identify
a paper stating precisely these contract formulas, but failure to locate
such a statement does not establish originality.

The next research question is whether their combination with canonical
source identity enables a useful new contract compiler, a nontrivial tractable
class under semantic constraints, or an error/storage optimization that existing
tools do not already provide. Until that is demonstrated, this document should
be treated as an internal technical note and a possible extension of the
existing relations manuscript, not a submission-ready standalone contribution.

\begin{thebibliography}{9}
\bibitem{gatterbauerSuciu}
Wolfgang Gatterbauer and Dan Suciu.
Oblivious bounds on the probability of Boolean functions.
\emph{ACM Transactions on Database Systems}, 39(1), Article 5, 2014.
\href{https://doi.org/10.1145/2532641}{doi:10.1145/2532641}.

\bibitem{kanagalDeshpande}
Bhargav Kanagal and Amol Deshpande.
Lineage processing over correlated probabilistic databases.
\emph{SIGMOD}, 2010.
\href{https://research.google/pubs/lineage-processing-on-correlated-probabilistic-databases/}{Author-hosted publication page}.

\bibitem{gopalan}
Parikshit Gopalan, Noam Nisan, Rocco A. Servedio, Kunal Talwar, and Avi Wigderson.
Smooth Boolean functions are easy: Efficient algorithms for low-sensitivity
functions. \emph{ITCS}, pages 59--70, 2016.
\href{https://doi.org/10.1145/2840728.2840738}{doi:10.1145/2840728.2840738}.

\bibitem{blanc}
Guy Blanc, Caleb Koch, Jane Lange, and Li-Yang Tan.
A query-optimal algorithm for finding counterfactuals.
\emph{Proceedings of Machine Learning Research}, 162:2075--2090, 2022.
\url{https://proceedings.mlr.press/v162/blanc22a.html}.

\bibitem{amarilliMonet}
Antoine Amarilli and Mika\"el Monet.
Weighted counting of matchings in unbounded-treewidth graph families.
\emph{MFCS}, 2022. Extended version, arXiv:2205.00851v2, 2023, Section 6.
\url{https://arxiv.org/abs/2205.00851}.

\bibitem{caiLuXia}
Jin-Yi Cai, Pinyan Lu, and Mingji Xia.
Holographic reduction, interpolation and hardness.
\emph{Computational Complexity}, 21(4):573--604, 2012.
\url{https://pages.cs.wisc.edu/~jyc/880/HRIH.pdf}.
\end{thebibliography}
\end{document}
